\documentclass[11pt]{article}
\usepackage{amsmath,amsfonts,amssymb,amsthm}
\usepackage{graphicx}
\usepackage{algorithm}
\usepackage{algorithmic}
\usepackage{float}
\usepackage{hyperref}
\usepackage{xcolor}

\newtheorem{definition}{Definition}
\newtheorem{remark}{Remark}

\title{Extreme learning machine-assisted solution of biharmonic equations via its coupled schemes}
\author{Xi'an Li, Jinran Wu, Zhe Ding, Xin Tai, Liang Liu, You-Gan Wang}
\date{}

\begin{document}
\maketitle

\section{Introduction}

Biharmonic equations are fourth-order partial differential equations that appear frequently in physics and applied mathematics, particularly in elasticity, Stokes flow problems, and thin plate splines. The standard biharmonic equation takes the form:
\begin{align}
\Delta^2 u(x) &= f(x), \quad \text{in } \Omega, \\
u(x) &= g(x), \quad \text{on } \partial\Omega, \\
\Delta u(x) &= h(x), \quad \text{on } \partial\Omega,
\end{align}
where $\Omega$ is a polygonal or polyhedral domain in $d$-dimensional Euclidean space with a piecewise Lipschitz boundary, $f(x) \in L^2(\Omega)$ is a prescribed function, and $\Delta$ represents the standard Laplace operator.

Traditional numerical methods for solving biharmonic equations include finite difference, finite volume, and finite element methods. However, these approaches face significant challenges when dealing with irregular domains and high-dimensional problems—often referred to as the "curse of dimensionality."

This paper introduces the Coupled Extreme Learning Machine (CELM) method, which combines the computational efficiency of Extreme Learning Machines (ELMs) with a coupled scheme approach to solve biharmonic equations.

\section{Methodology}

\subsection{Extreme Learning Machine (ELM)}

The Extreme Learning Machine is a single hidden layer fully connected neural network designed to map from $\mathbb{R}^d$ to $\mathbb{R}$. The distinctive feature of ELM is that the input layer weights and biases are randomly assigned and fixed during training, with only the output weights analytically determined.

Mathematically, the ELM mapping $\ell(x)$ with an input $x \in \mathbb{R}^d$ is expressed as:
\begin{equation}
\ell(x) = \sum_{i=1}^{m} \beta_i \sigma_i(W^{[i]}x + b_i),
\end{equation}
where:
\begin{itemize}
    \item $W = [W^{[1]}, W^{[2]}, \ldots, W^{[m]}]^T \in \mathbb{R}^{m \times d}$ with $W^{[i]} \in \mathbb{R}^d$ being the weight of the $i$-th hidden unit
    \item $b = (b_1, b_2, \ldots, b_m)^T \in \mathbb{R}^m$ is the bias vector
    \item $\sigma(\cdot)$ is an element-wise non-linear activation function
    \item $\beta = [\beta_1, \beta_2, \ldots, \beta_m]^T \in \mathbb{R}^m$ is the weight vector connecting the hidden units to the output
\end{itemize}

\subsection{Coupled Extreme Learning Machine (CELM) for Biharmonic Equations}

The CELM approach transforms the original fourth-order biharmonic equation into two coupled second-order Poisson equations by introducing an auxiliary variable $v = \Delta u$:
\begin{align}
\begin{cases}
\Delta v = f, & \text{in } \Omega \\
v = h, & \text{on } \partial\Omega
\end{cases}
\quad \text{and} \quad
\begin{cases}
\Delta u = v, & \text{in } \Omega \\
u = g, & \text{on } \partial\Omega
\end{cases}
\end{align}

The CELM method uses two ELM networks to approximate the solutions $u$ and $v$:
\begin{align}
u(x; \tilde{W}, \tilde{b}) &= \sigma(x^T \cdot \tilde{W} + \tilde{b}^T) \cdot \tilde{\beta} \\
v(x; \hat{W}, \hat{b}) &= \sigma(x^T \cdot \hat{W} + \hat{b}^T) \cdot \hat{\beta}
\end{align}

Substituting the approximation for $v$ into the first equation yields:
\begin{align}
\Delta \left( \sum_{i=1}^{K} \hat{\beta}_i \sigma(V_i(x)) \right) &= f(x) \\
\sum_{i=1}^{K} \hat{\beta}_i \sigma(V_i(x)) &= h(x)
\end{align}
where $V_i(x) = x^T \cdot \hat{W}^{[i]} + \hat{b}_i$ is the linear output for the $i$-th hidden unit of the ELM, and $K$ is the number of hidden units.

Similarly, for the second equation:
\begin{align}
\Delta \left( \sum_{j=1}^{N} \tilde{\beta}_j \sigma(U_j(x)) \right) &= \sum_{i=1}^{K} \hat{\beta}_i \sigma(V_i(x)) \\
\sum_{j=1}^{N} \tilde{\beta}_j \sigma(U_j(x)) &= g(x)
\end{align}
where $U_j(x) = x^T \cdot \tilde{W}^{[j]} + \tilde{b}_j$, and $N$ is the number of hidden units.

\subsection{Algorithm Implementation}

The CELM algorithm involves the following steps:

\begin{enumerate}
    \item Generating collocation points: Sample $Q$ random points $\{x_I^q\}_{q=1}^Q$ from the interior of the domain $\Omega$ and $P$ random points $\{x_B^p\}_{p=1}^P$ from the boundary $\partial\Omega$.
    
    \item Computing residuals at the collocation points:
    \begin{align}
    R_v(x_I^q) &= \Delta \left( \sum_{i=1}^{K} \hat{\beta}_i \sigma(V_i(x_I^q)) \right) - f(x_I^q) \\
    R_u(x_I^q) &= \Delta \left( \sum_{j=1}^{N} \tilde{\beta}_j \sigma(U_j(x_I^q)) \right) - \sum_{i=1}^{K} \hat{\beta}_i \sigma(V_i(x_I^q)) \\
    L_v(x_B^p) &= \sum_{i=1}^{K} \hat{\beta}_i \sigma(V_i(x_B^p)) - h(x_B^p) \\
    L_u(x_B^p) &= \sum_{j=1}^{N} \tilde{\beta}_j \sigma(U_j(x_B^p)) - g(x_B^p)
    \end{align}
    
    \item Setting up the linear system $H\beta = S$ where:
    \begin{equation}
    H = 
    \begin{bmatrix}
    O & H_{12} \\
    H_{21} & -H_{22} \\
    H_{31} & O \\
    O & H_{42}
    \end{bmatrix},
    \quad
    \beta = 
    \begin{bmatrix}
    \tilde{\beta} \\
    \hat{\beta}
    \end{bmatrix},
    \quad
    S = 
    \begin{bmatrix}
    S_f \\
    S_0 \\
    S_g \\
    S_h
    \end{bmatrix}
    \end{equation}
    
    with matrix blocks:
    \begin{itemize}
        \item $H_{12}$ contains second derivatives of activation functions evaluated at interior points: $\sigma''(V_i(x_I^q))$
        \item $H_{21}$ contains second derivatives of activation functions: $\sigma''(U_j(x_I^q))$
        \item $H_{22}$ contains activation function values: $\sigma(V_i(x_I^q))$
        \item $H_{31}$ contains activation function values at boundary points: $\sigma(U_j(x_B^p))$
        \item $H_{42}$ contains activation function values at boundary points: $\sigma(V_i(x_B^p))$
    \end{itemize}
    
    and right-hand side vectors:
    \begin{itemize}
        \item $S_f = [f(x_I^1), f(x_I^2), \ldots, f(x_I^Q)]^T$
        \item $S_0 = [0, 0, \ldots, 0]^T$ (zero vector of length $Q$)
        \item $S_g = [g(x_B^1), g(x_B^2), \ldots, g(x_B^P)]^T$
        \item $S_h = [h(x_B^1), h(x_B^2), \ldots, h(x_B^P)]^T$
    \end{itemize}
    
    \item Solving the linear system for $\beta$ using one of the following approaches:
    \begin{itemize}
        \item If $H$ is square and invertible: $\beta = H^{-1}S$
        \item If $H$ is rectangular and full rank: $\beta = H^\dagger S$ where $H^\dagger = (H^T H)^{-1}H^T$ is the pseudo-inverse
        \item If $H$ is singular: Use Tikhonov regularization with $\beta = (H^T H + \lambda I)^{-1}HS$ where $\lambda > 0$ is the ridge parameter
    \end{itemize}
\end{enumerate}

\subsection{Activation Functions and Parameter Initialization}

The paper investigates the influence of different activation functions on the CELM performance:
\begin{itemize}
    \item Hyperbolic tangent (tanh)
    \item Gaussian function
    \item Sine function
    \item Trigonometric combination (sine + cosine)
\end{itemize}

\begin{remark}[Lipschitz Continuity]
If an activation function $\sigma$ is continuous (i.e., $\sigma \in C^1$) and satisfies the boundedness condition:
\begin{align}
|\sigma(x)| < 1 \quad \text{and} \quad |\sigma'(x)| < 1
\end{align}
for any $x \in \mathbb{R}$, then:
\begin{align}
|\sigma(x) - \sigma(y)| \leq |x-y| \quad \text{and} \quad |\sigma'(x) - \sigma'(y)| < |x-y|
\end{align}
for any $x, y \in \mathbb{R}$. All the activation functions mentioned above satisfy this condition, ensuring boundedness and robust regularity.
\end{remark}

For parameter initialization, the paper explores several approaches:
\begin{itemize}
    \item Normal distribution initialization (mean=0, standard deviation=10)
    \item Unitary uniform distribution initialization in $[-1, 1]$
    \item Scale uniform distribution initialization in $[-\sigma, \sigma]$ with scale factor $\sigma = 10$
    \item Mixed initialization (normal for weights, scale uniform for biases)
\end{itemize}

\section{Numerical Examples and Results}

The paper evaluates the CELM method using four numerical examples with different domain complexities and dimensionalities.

\subsection{Performance Metrics}

Two error metrics are used to assess the CELM performance:
\begin{itemize}
    \item Point-wise absolute error: $APE(x_i) = |\tilde{u}(x_i) - u^*(x_i)|$
    \item Relative error: $REL = \sqrt{\frac{\sum_{i=1}^{N'} |\tilde{u}(x_i) - u^*(x_i)|^2}{\sum_{i=1}^{N'} |u^*(x_i)|^2}}$
\end{itemize}
where $\tilde{u}(x_i)$ is the CELM approximation, $u^*(x_i)$ is the exact solution, and $N'$ is the number of test points.

\subsection{Example 1: Regular 2D Domain}

The first example solves the biharmonic equation in a square domain $\Omega = [-1, 1] \times [-1, 1]$ with the analytical solution $u(x_1, x_2) = \sin(\pi x_1)\sin(\pi x_2)$ and corresponding force term $f(x_1, x_2) = 4\pi^4 \sin(\pi x_1)\sin(\pi x_2)$. Boundary conditions are $g(x_1, x_2) = 0$ and $h(x_1, x_2) = 0$ on $\partial\Omega$.

Results:
\begin{itemize}
    \item CELM with sine and sine+cosine activation functions achieved the lowest errors (around $10^{-14}$) with scale uniform initialization.
    \item Hyperbolic tangent and Gaussian activation functions produced higher errors ($10^{-6}$ to $10^{-9}$).
    \item The relative error did not decrease with increasing numbers of hidden units (beyond a threshold), indicating stable performance.
    \item Runtime increased superlinearly with more hidden neurons, but the increase was gradual.
    \item The CELM model remained stable and robust across different scale factors $\sigma$.
    \item Comparison with the finite difference method (FDM) showed CELM significantly outperformed FDM in both accuracy and runtime, especially for dense grids.
\end{itemize}

\subsection{Example 2: Irregular 2D Domain}

The second example addresses the biharmonic equation in a hexagram domain derived from the unit square $[0,1] \times [0,1]$. The analytical solution is $u(x_1, x_2) = 10(x_1 - 2x_1^3 + x_1^4)(x_2 - 2x_2^3 + x_2^4)$.

Results:
\begin{itemize}
    \item On this irregular domain, CELM with sine activation function and scale uniform initialization achieved the best performance with relative errors around $10^{-10}$.
    \item The model showed excellent stability and robustness with varying numbers of hidden nodes and scale factors.
    \item The coefficient matrix remained full rank, eliminating the need for Tikhonov regularization.
\end{itemize}

\subsection{Example 3: 3D Domain with Complex Geometry}

The third example tackles the biharmonic equation in a three-dimensional cubic domain $\Omega = [0,1] \times [0,1] \times [0,1]$ with one large hole and eight smaller holes. The analytical solution is $u(x_1, x_2, x_3) = 10(x_1 - 2x_1^3 + x_1^4)(x_2 - 2x_2^3 + x_2^4)(x_3 - 2x_3^3 + x_3^4)$.

Results:
\begin{itemize}
    \item CELM with sine and sine+cosine activation functions demonstrated superior performance (relative errors around $10^{-13}$) compared to hyperbolic tangent and Gaussian functions.
    \item Scale uniform initialization provided the best results for parameter initialization.
    \item The method showed favorable stability and robustness with increasing numbers of hidden nodes and varying scale factors.
\end{itemize}

\subsection{Example 4: 3D Spherical Shell Domain}

The final example examines the biharmonic equation on a closed domain $\Omega$ bounded by two spherical surfaces with radii $r_1 = 0.5$ and $r_2 = 1.0$. The exact solution is $u(x_1, x_2, x_3) = \sin(\pi x_1)\sin(\pi x_2)\sin(\pi x_3)$.

Results:
\begin{itemize}
    \item For this spherical shell geometry, CELM with sine and sine+cosine activation functions again achieved the lowest errors (around $10^{-9}$) when using scale uniform initialization.
    \item The point-wise error distribution confirmed that these activation functions yielded significantly lower errors than hyperbolic tangent and Gaussian functions.
    \item The method maintained stability and robustness across varying numbers of hidden nodes and scale factors.
\end{itemize}

\section{Implications and Advantages}

The CELM method offers several significant advantages for solving biharmonic equations:

\begin{enumerate}
    \item \textbf{Mesh-Free Approach}: Unlike traditional numerical methods like finite elements or finite differences, CELM operates without requiring mesh generation, which is particularly valuable for complex geometries.
    
    \item \textbf{Computational Efficiency}: By avoiding the iterative gradient descent-based techniques used in physics-informed neural networks (PINNs), CELM provides a more efficient solution method.
    
    \item \textbf{High Accuracy}: The numerical examples demonstrate that CELM achieves remarkably high accuracy, often with relative errors in the range of $10^{-9}$ to $10^{-14}$, which outperforms traditional methods.
    
    \item \textbf{Stability and Robustness}: The method shows stable performance across varying numbers of hidden units and different scale factors for initialization, indicating robustness.
    
    \item \textbf{Flexibility for Complex Geometries}: CELM effectively handles both regular and irregular domains in two and three dimensions, including geometries with holes and spherical shells.
    
    \item \textbf{Activation Function Selection}: The study reveals that sine and trigonometric (sine+cosine) activation functions significantly outperform hyperbolic tangent and Gaussian functions for biharmonic problems.
    
    \item \textbf{Initialization Strategy Impact}: Scale uniform initialization of weights and biases consistently produces the best results, providing practical guidance for implementing the method.
\end{enumerate}

The CELM method shows potential for extension to other types of PDEs beyond biharmonic equations. The same approach could be applied to various steady-state and time-dependent scenarios, as well as inverse problems where traditional methods face significant challenges.

\end{document} 